% a4-work.tex — preamble for the a4-work pandoc layout.
% Loaded via `include-in-header` from a4-work.yaml (auto-discovered by 2pdf.sh).
% Formal A4 output for sharing with WDODelta colleagues and management.
%
% Styled to the WDODelta huisstijl (official house style, June 2021):
% donkerblauw #075895 for headings/links/footer, blauw #00b0ea as a single
% restrained accent (title rule only), Carlito as the Calibri-compatible
% fallback font (Vivala Sans Rounded is proprietary and unavailable here),
% near-black body text (#1a1a1a), footer contact line. No logo — the huisstijl
% doc marks the raster preview asset "niet voor productie"; only the official
% SVG (request via Afdeling Communicatie) belongs in production documents.

\usepackage{xcolor}
\definecolor{wdoddark}{HTML}{075895}
\definecolor{wdodaccent}{HTML}{00b0ea}
\definecolor{wdodtext}{HTML}{1a1a1a}

% Carlito is missing a glyph for U+2011 (non-breaking hyphen), which shows up
% in Dutch compound words (e.g. "Epic‑bord", "T‑shirt‑sizing") and rendered as
% a boxed replacement character. Map it to a regular hyphen wrapped in \mbox
% so it still doesn't break across a line — same visual + behavior Jelle asked
% for ("that should be a -"). Add further \newunicodechar lines here if other
% source documents surface more Carlito glyph gaps.
\usepackage{newunicodechar}
\newunicodechar{‑}{\mbox{-}}

\usepackage{titlesec}

% Heading hierarchy with 11pt body — v2 (2026-07-28, widened on Jelle's feedback:
% "font size differences between headings should be larger, level 1/2/3/4 not
% clearly distinguishable"). Explicit \fontsize{size}{baselineskip} instead of
% LaTeX's built-in \LARGE/\Large/\large steps, whose default increments are too
% small (17.28/14.4/12.0pt) to read as a confident staircase.
%   H1 (\section)       — 26pt, bold, wdoddark, accent rule beneath
%   H2 (\subsection)    — 19pt, bold, wdoddark
%   H3 (\subsubsection) — 15pt, bold, wdoddark
%   H4 (\paragraph)     — 12.5pt, bold, wdoddark, own block line (not run-in)
%   Body                — 11.0pt, wdodtext (near-black, not pure black)
%
% Gaps: H1→H2 = 7pt (37%), H2→H3 = 4pt (27%), H3→H4 = 2.5pt (20%), H4→body = 1.5pt (14%)
% Each step widens going up — mirrors how a reader scans top-down for structure.
% All bold — weight adds to size to make each level unmistakable at a glance.
% Restraint per huisstijl ("beperk het aantal huisstijlkleuren per scherm"):
% only wdoddark carries headings/links; wdodaccent is reserved for the H1 rule
% and the callout box's left bar (wdocallout, below) — no other use of it.

\titleformat{\section}
  {\normalfont\fontsize{26}{31}\selectfont\bfseries\color{wdoddark}}{\thesection}{0.6em}{}[\vspace{2pt}{\color{wdodaccent}\titlerule[1.2pt]}]

\titleformat{\subsection}
  {\normalfont\fontsize{19}{23}\selectfont\bfseries\color{wdoddark}}{\thesubsection}{0.6em}{}

\titleformat{\subsubsection}
  {\normalfont\fontsize{15}{18}\selectfont\bfseries\color{wdoddark}}{\thesubsubsection}{0.6em}{}

\titleformat{\paragraph}
  {\normalfont\fontsize{12.5}{15}\selectfont\bfseries\color{wdoddark}}{\theparagraph}{0.6em}{}

% Vertical spacing — generous above headings, tight below. Scaled up with the
% larger font sizes so whitespace stays proportional to the new staircase.
% Format: {left}{space above}{space below}
\titlespacing*{\section}
  {0pt}{2.6ex plus 0.6ex minus 0.4ex}{1.0ex plus 0.2ex}
\titlespacing*{\subsection}
  {0pt}{2.0ex plus 0.5ex minus 0.3ex}{0.8ex plus 0.2ex}
\titlespacing*{\subsubsection}
  {0pt}{1.6ex plus 0.4ex minus 0.2ex}{0.6ex plus 0.2ex}
\titlespacing*{\paragraph}
  {0pt}{1.3ex plus 0.3ex minus 0.2ex}{0.5ex plus 0.1ex}

% Body text: near-black per huisstijl (#1a1a1a), not pure black.
\AtBeginDocument{\color{wdodtext}}

% Callout box — target of lua/a4-work/callouts.lua (Obsidian callouts).
% One sober style for every callout type, per huisstijl restraint ("beperk
% het aantal huisstijlkleuren per scherm"): a very light tint of the accent
% blue as background arcering plus a 3pt accent bar on the left. breakable
% so a callout can cross a page boundary without clipping.
\usepackage{tcolorbox}
\tcbuselibrary{breakable,skins}
\definecolor{wdodcalloutbg}{HTML}{EAF6FC}

% Shared skeleton for both box variants below — breakable/enhanced/frameless
% plumbing is identical; only colour, bar weight, corner, and padding differ.
\tcbset{wdocbase/.style={breakable, enhanced, frame hidden, boxrule=0pt}}

\newtcolorbox{wdocallout}{
  wdocbase, arc=1.5pt,
  colback=wdodcalloutbg,
  borderline west={3pt}{0pt}{wdodaccent},
  left=10pt, right=10pt, top=7pt, bottom=7pt,
  before skip=10pt, after skip=10pt
}

% Plain blockquotes ("> tekst" with no "[!type]" marker) fall straight through
% the Lua filter above untouched — it only boxes real Obsidian callouts. Left
% at LaTeX's built-in `quote` environment, that renders as a bare indent with
% no frame at all ("indent maar geen kader", Jelle 2026-08-18). Obsidian's own
% editor gives every blockquote a left bar regardless of a [!type] marker, so
% readers expect one here too. Quiet, uncoloured bar — lighter than
% wdocallout's tinted box — keeps the two visually distinct: a callout is a
% highlighted note, a plain quote is a quiet aside.
\newtcolorbox{wdocquote}{
  wdocbase, arc=0pt,
  colback=white,
  borderline west={2pt}{0pt}{wdoddark!45},
  left=10pt, right=6pt, top=5pt, bottom=5pt,
  before skip=8pt, after skip=8pt
}
\renewenvironment{quote}{\begin{wdocquote}}{\end{wdocquote}}

% Footer — org name + page number, satisfying the huisstijl Self-Check
% "footer/contact reference" rule. No header (no logo asset to place there).
\usepackage{fancyhdr}
\pagestyle{fancy}
\fancyhf{}
\renewcommand{\headrulewidth}{0pt}
\renewcommand{\footrulewidth}{0.4pt}
\fancyfoot[L]{\small\color{wdoddark}Waterschap Drents Overijsselse Delta (WDODelta)}
\fancyfoot[R]{\small\color{wdoddark}\thepage}
